\paraid{p-c05f83eda3c9}
We represent a compact metric space by its distance functions in
an infinite-dimensional function space.
Using the continuous invariant probabilities of Part~\ref{part:measures},
we construct finite-dimensional approximations of this representation
that are equivariant under isometries and vary continuously with the compact space.
Let
$\GHSpace$
denote the space of isometry classes of nonempty compact metric spaces.
Our main result is \Cref{thm:local-models}.

\paraid{p-79876a468ed5}
The symmetry of the input space is part of the approximation problem.
Rouyer's results
\cite[arXiv v1, Theorems~2 and~4]{Rouyer2011}
imply that compact metric spaces with trivial isometry group are dense in
$\GHSpace$.
The approximation must also remain continuous at spaces with nontrivial
isometries.
We retain whole finite spectral subspaces and allow changes of their
orthonormal bases.
The resulting orthogonal actions express both the symmetry of the input
and the coordinate changes needed when the input varies.


\paraid{p-a0f40ab05c70}
For every nonempty compact metric space
$\MetricSpace{\BaseCarrier}{\MetricSymbol}$
and every
$\LocalTolerance>0$,
we construct a neighborhood
$\ModelNeighborhood\subset\GHSpace$
and a fixed integer
$\ModelDimension\geq1$.
For each compact metric space
$\MetricSpace{\ComparisonCarrier}{\ComparisonMetric}$
whose isometry class belongs to
$\ModelNeighborhood$,
there is a norm
$\NormSymbol_\ComparisonCarrier$
on
$\RealNumbers^\ModelDimension$
and a continuous map
\[
 \FeatureCoordinates_\ComparisonCarrier\colon\ComparisonCarrier
 \to\RealNumbers^\ModelDimension.
\]
For
$\ComparisonPoint,\IntegrationPoint\in\ComparisonCarrier$,
define the induced pseudometric by
\[
 \Pseudometric_\ComparisonCarrier(\ComparisonPoint,\IntegrationPoint)
 =\NormSymbol_\ComparisonCarrier
 (\FeatureCoordinates_\ComparisonCarrier(\ComparisonPoint)
 -\FeatureCoordinates_\ComparisonCarrier(\IntegrationPoint)).
\]
Its uniform error is
\[
 \LocalError\MetricSpace{\ComparisonCarrier}{\ComparisonMetric}
 =\sup_{\ComparisonPoint,\IntegrationPoint\in\ComparisonCarrier}
 |\Pseudometric_\ComparisonCarrier(\ComparisonPoint,\IntegrationPoint)
 -\ComparisonMetric(\ComparisonPoint,\IntegrationPoint)|.
\]
The function
$\LocalError$
is continuous on
$\ModelNeighborhood$
and satisfies
$\LocalError\MetricSpace{\BaseCarrier}{\MetricSymbol}<\LocalTolerance$.
We can
shrink
$\ModelNeighborhood$
so that
$\LocalError\MetricSpace{\ComparisonCarrier}{\ComparisonMetric}<\LocalTolerance$
for every
$\MetricSpace{\ComparisonCarrier}{\ComparisonMetric}\in\ModelNeighborhood$.
The dimension and the neighborhood may depend on the center
$\MetricSpace{\BaseCarrier}{\MetricSymbol}$
and the tolerance
$\LocalTolerance$.
Within that neighborhood,
the image
$\FeatureCoordinates_\ComparisonCarrier(\ComparisonCarrier)$
is a compact subset of one fixed vector space equipped with the varying norm
$\NormSymbol_\ComparisonCarrier$.
We also use
$\NormSymbol_\ComparisonCarrier$
to denote the metric
$(\FirstVector,\SecondVector)\mapsto\NormSymbol_\ComparisonCarrier(\FirstVector-\SecondVector)$
on
$\RealNumbers^\ModelDimension$.
The correspondence
$\{(\ComparisonPoint,\FeatureCoordinates_\ComparisonCarrier(\ComparisonPoint))
 \mid\ComparisonPoint\in\ComparisonCarrier\}$
has distortion at most
$\LocalError\MetricSpace{\ComparisonCarrier}{\ComparisonMetric}$.
Then \Cref{thm:gh-correspondence} implies
\[
 \GHDistance\!\left(
 \MetricSpace{\ComparisonCarrier}{\ComparisonMetric},
 \MetricSpace{\FeatureCoordinates_\ComparisonCarrier(\ComparisonCarrier)}{\NormSymbol_\ComparisonCarrier}
 \right)<\frac{\LocalTolerance}{2}.
\]
We summarize the continuity assertion of \Cref{thm:local-models}\ref{item:model-continuity} for the norms and point maps.
Let
$\{\MetricSpace{\ComparisonCarrier_\SequenceIndex}{\ComparisonMetric_\SequenceIndex}\}_{\SequenceIndex\in\NonnegativeIntegers}$
be a sequence in
$\ModelNeighborhood$,
and let
$\MetricSpace{\ComparisonCarrier}{\ComparisonMetric}\in\ModelNeighborhood$.
Assume that these spaces are embedded isometrically in a common compact ambient space
and that their carriers converge in Hausdorff distance to
$\ComparisonCarrier$.
After fixing these spaces and embeddings, choose a coordinate pair
$(\FeatureCoordinates_\ComparisonCarrier,\NormSymbol_\ComparisonCarrier)$
for the limit space and coordinate pairs
$\{(\FeatureCoordinates_{\ComparisonCarrier_\SequenceIndex},\NormSymbol_{\ComparisonCarrier_\SequenceIndex})\}_{\SequenceIndex\in\NonnegativeIntegers}$
for the spaces in the sequence so that
\[
 \sup_{\substack{\CoefficientVector\in\RealNumbers^\ModelDimension\\\norm{\CoefficientVector}_2=1}}
 \bigl|\NormSymbol_{\ComparisonCarrier_\SequenceIndex}(\CoefficientVector)
       -\NormSymbol_\ComparisonCarrier(\CoefficientVector)\bigr|\to0,
\]
and for every sequence
$\ComparisonPoint_\SequenceIndex\in\ComparisonCarrier_\SequenceIndex$
and every
$\ComparisonPoint\in\ComparisonCarrier$
with
$\ComparisonPoint_\SequenceIndex\to\ComparisonPoint$
in the common ambient space,
the coordinate values satisfy
\[
 \FeatureCoordinates_{\ComparisonCarrier_\SequenceIndex}(\ComparisonPoint_\SequenceIndex)
 \to\FeatureCoordinates_\ComparisonCarrier(\ComparisonPoint).
\]
For each
$\OrthogonalChange\in\OrthogonalGroup(\ModelDimension)$,
the corresponding change of orthonormal basis acts on the pair by
\[
 (\FeatureCoordinates_\ComparisonCarrier,\NormSymbol_\ComparisonCarrier)
 \longmapsto
 (\OrthogonalChange\circ\FeatureCoordinates_\ComparisonCarrier,
  \NormSymbol_\ComparisonCarrier\circ\OrthogonalChange^{-1}).
\]
This action preserves the induced pseudometric.

\paraid{p-210d146c099f}
The feature map is equivariant with respect to an orthogonal representation
of the isometry group of each input space.
In Remark~\ref{rem:local-isometry-representation},
we construct a continuous
homomorphism
\[
 \IsometryRepresentation{\ComparisonCarrier}{\ComparisonMetric}
 \colon\Isom\MetricSpace{\ComparisonCarrier}{\ComparisonMetric}
 \to\OrthogonalGroup(\ModelDimension)
\]
that satisfies equivariance and preserves the chosen norm.
For every
$\FirstIsometry\in\Isom\MetricSpace{\ComparisonCarrier}{\ComparisonMetric}$
and every
$\ComparisonPoint\in\ComparisonCarrier$,
the equivariance identity is
\[
 \FeatureCoordinates_\ComparisonCarrier(\FirstIsometry\ComparisonPoint)
 =\IsometryRepresentation{\ComparisonCarrier}{\ComparisonMetric}(\FirstIsometry)
   \FeatureCoordinates_\ComparisonCarrier(\ComparisonPoint).
\]
For every
$\FirstIsometry\in\Isom\MetricSpace{\ComparisonCarrier}{\ComparisonMetric}$
and every
$\CoefficientVector\in\RealNumbers^\ModelDimension$,
the norm-preservation identity is
\[
 \NormSymbol_\ComparisonCarrier
   \bigl(\IsometryRepresentation{\ComparisonCarrier}{\ComparisonMetric}(\FirstIsometry)
          \CoefficientVector\bigr)
 =\NormSymbol_\ComparisonCarrier(\CoefficientVector).
\]
Thus the varying isometry groups act through one fixed orthogonal group
on the model neighborhood.
Orthogonal changes of basis remain relevant even when the input space
has trivial isometry group.

\paraid{p-852d0b9cfe26}
We first recall two examples of spectral methods in geometry that provide context for our construction.
Bates
\cite[arXiv v1, Theorem~2]{Bates2014Eigenfunctions}
proves that closed connected Riemannian manifolds of fixed dimension at least two
and unit volume,
with a common lower bound on Ricci curvature and a common positive lower bound
on injectivity radius,
admit smooth embeddings by a uniformly bounded number of Laplacian eigenfunctions.
Kuwae and Shioya study spectral structures associated with nonnegative
self-adjoint operators on varying Hilbert spaces.
Under compact convergence of these structures,
they prove eventual equality of the dimensions of spectral subspaces
for bounded intervals whose endpoints lie outside the limiting spectrum
\cite[Theorem~2.6]{KuwaeShioya2003Spectral}.
In the present  paper,
 we use the distance kernel on arbitrary compact metric spaces
and approximate the distance uniformly.
We use the continuous invariant full-support assignment of probabilities
constructed in \Cref{thm:invariant-measures} of Part~\ref{part:measures}.
For
$\MetricSpace{\BaseCarrier}{\MetricSymbol}$,
the associated distance operator
$\DistanceOperator{\BaseCarrier}{\MetricSymbol}
 \colon\LebesgueSpace^2(\BaseCarrier,\SelectedMeasure{\BaseCarrier}{\MetricSymbol})
 \to\LebesgueSpace^2(\BaseCarrier,\SelectedMeasure{\BaseCarrier}{\MetricSymbol})$
is a self-adjoint Hilbert--Schmidt integral operator.
For every
$f\in\LebesgueSpace^2(\BaseCarrier,\SelectedMeasure{\BaseCarrier}{\MetricSymbol})$
and
$\BasePoint\in\BaseCarrier$,
its value is
\[
 (\DistanceOperator{\BaseCarrier}{\MetricSymbol}f)(\BasePoint)
 =\int_{\BaseCarrier}\MetricSymbol(\BasePoint,\IntegrationPoint)
 f(\IntegrationPoint)\,d\SelectedMeasure{\BaseCarrier}{\MetricSymbol}(\IntegrationPoint).
\]
Maria,
Oudot,
and Solomon study this operator and its spectral coordinates
\cite[Section~3]{MariaOudotSolomon2020}.
We approximate the distance functions by unique best approximations in
finite spectral subspaces,
using
$\LebesgueSpace^\IntegrabilityExponent$
norms.
Here
$2\leq\IntegrabilityExponent<\infty$,
and the norm is taken with respect to the selected probability measure.
We choose the exponent so that the distances between the approximants
approximate the original distances uniformly.
We retain whole eigenspaces and account for eigenvalue multiplicities
by orthogonal changes of basis.
Using continuity of the spectral subspaces and of the best approximations,
we prove the asserted continuity of the norms and point maps.

% Retired paragraph ID from the introduction to former Subsections 8.5 and 8.6.
% \paraid{p-fc6db22afc8b}

\paraid{p-cfb0b7538b10}
By the interpolation result in
\cite[Theorem~1.1]{Ishiki2026Interpolation},
we obtain a uniform error bound when extending prescribed metrics from a
discrete family of closed subsets of a fixed metrizable space.
Here we use \Cref{thm:local-models} to construct maps and norms on varying
compact spaces for the subsequent absolute retract argument.

\medskip
\noindent\textbf{Dependence on the other parts.}
\paraid{p-b4ac06e8df4d}
\Cref{thm:invariant-measures} is the input from Part~\ref{part:measures}.
We prove the spectral approximation and its applications using this
assignment and the preliminaries of Section~\ref{sec:prelim-spectral}.
Part~\ref{part:topology} uses \Cref{thm:local-models} to represent the varying
norms in fixed Banach spaces and combine the local point maps.
The orthogonal actions determine the compact group used in that construction.

\medskip
\noindent\textbf{Organization.}
\paraid{p-6ad9276d62d7}
In Section~\ref{sec:prelim-spectral},
we recall best approximation in
$\LebesgueSpace^\IntegrabilityExponent$,
spectral decomposition and perturbation,
and square roots of positive definite matrices.
In Subsection~\ref{subsec:spectral-distance-operator},
we define the distance operator and study its finite
spectral subspaces.
In Subsection~\ref{subsec:spectral-distance-approximation},
we prove approximation of distance functions using finite
spectral subspaces and unique best approximations in finite
$\LebesgueSpace^\IntegrabilityExponent$
norms.
In Subsection~\ref{subsec:spectral-subspace-continuity},
we prove continuity of the finite spectral subspaces
and construct orthonormal bases whose continuous extensions converge uniformly
on a common compact ambient space,
as asserted in \ref{item:spectral-basis-uniform-convergence}
of \Cref{lem:spectral-continuity}.
In Subsection~\ref{subsec:spectral-local-models},
we combine these results to construct the local maps and
varying norms in \Cref{thm:local-models}.

\medskip
\noindent\textbf{Conventions and notation.}
\paraid{p-21fe15dbd1cc}
All compact metric spaces are nonempty,
and finite-dimensional coordinate spaces are real.
Weak convergence of probabilities and the uniform integral estimates are
recalled in Subsection~\ref{subsec:preliminaries-probability}.
Function spaces are real unless complex scalars are explicitly indicated
for spectral perturbation.
An isometry is a surjective isometric embedding.
The pair
$\MetricSpace{\BaseCarrier}{\MetricSymbol}$
specifies both the underlying set and its metric.
In particular,
$\SelectedMeasure{\BaseCarrier}{\MetricSymbol}$
and
$\DistanceOperator{\BaseCarrier}{\MetricSymbol}$
depend on that pair.
When compact spaces are embedded isometrically into a common ambient
space,
we use the same symbols for the embedded subsets and the restricted
metrics.
The notation
$\IdentityMap_\BaseCarrier$
denotes the identity map on
$\BaseCarrier$.
We omit the subscript when its domain is specified.
Sequences are indexed by
$\NonnegativeIntegers=\{0,1,2,\ldots\}$.
